\section{Early source curvature makes the finite owner bound linear}
Use the actual arms and mixed/source-prefix premises (1) of
\texttt{arbitrary\_source\_gap\_finite\_owner.tex}.
Retain $D=a-\alpha-h$, $d=\kappa_Y-\alpha$ and
$T=e_{n+1}-e_n$. No fixed upper bound on these parameters is imposed.

\begin{theorem}[Linear finite-owner bound at arbitrary source gap]
If $T>0$, then
\[
 \boxed{d\ge6,\qquad D\le\min\{2d(d-3),\,12d-25\},
       \qquad a\le\alpha+24d-50.}                    \tag{1}
\]
In particular, at each fixed opposite first missing weight $\kappa$,
\[
 \boxed{a\le\min\{1+4(\kappa-1)(\kappa-4),\,24\kappa-73\},
       \qquad j\le a+1.}                             \tag{2}
\]
The constants are explicit sufficient bounds; no optimality is claimed.
For every mixed S edge still outside the previous count-mass theorem,
the same source-prefix premises therefore imply
\[
                 h\le12\kappa-37,\qquad g\le60\kappa-179.       \tag{3}
\]
Thus both counts, the owner weight and the shared forward rank are
bounded LINEARLY in the first missing weight. The opposite weight,
capacities and initial padding are not bounded by this statement.
\end{theorem}

\begin{proof}
The preceding theorem supplies $d\ge6$, the quadratic bound,
and the exact source/primitive identities. It suffices to show
that $D\ge12d-24$ forces $T\le0$.
The cases $r=a-j\le-2$ have no primitive curvature.
At $r=-1$, positive recovery requires $D+h+2-d<0$, which gives
$D\le d-5\le12d-25$ and contradicts the new degree assumption.

Suppose $r\ge0$ and $T>0$. As before set
\[
 e=D-2r-2\ge0,\qquad C=d_{r+1}-d_r>0,\qquad
 (r+1)C\le(e+1)d_r,
\]
where $d_i=(\Delta q_Z(i))^+$. If $s=d-e-2\le0$, the
FIRST source birth already pays. Otherwise put $L=s-1=d-e-3\ge0$.
The new degree assumption gives BOTH
\[
                  r\ge3e+6L+5,\qquad r\ge4e+3.       \tag{4}
\]
Indeed substitute $d=e+L+3$, $D=2r+2+e$; these are just
linear inequalities with $e,L\ge0$.

If $L=0$, the first birth is at least $q_Z(r)\ge d_r$, and
$r\ge e$ pays the curvature.
If $L=1$, the first birth is at least $q_Z(r-1)\ge d_{r-1}$.
The skeleton-only CI ratio gives
$r d_r\le(r+e+3)d_{r-1}$, while $r\ge2e+3$ ensures
$(e+1)(r+e+3)\le(r+1)r$.
Thus this case is also paid by the first birth.

Let $L\ge2$. The proof now uses a LATER compulsory source increment,
not a differentiated coefficient lower bound. The genuinely proved
early source-curvature profile gives
\[
       w_{\kappa+L}:=\Delta E_Y(\kappa+L)\ge2^{L-2},     \tag{5}
\]
because $L\le\kappa-3$, which follows from
$\kappa=\alpha+e+L+3$ and $\alpha\ge1$.
The term (5) is within the assumed nondecreasing source prefix:
its associated ordinary coefficient has index
\[
                 j+1-\kappa-L=r+h-2L\ge h.
\]
Thus (5) and the positive source convolution give
\[
 \Delta(E_YF_Z)_{j+1}\ge2^{L-2}F_Z(j+1-\kappa-L)
     =2^{L-2}q_Z(r+\kappa-1+L)
     \ge2^{L-2}q_Z(r-2L)
     \ge2^{L-2}d_{r-2L}.                              \tag{6}
\]
The reflected indices in the middle have endpoint sum EXACTLY
$\alpha+D=a-h$; both are nonnegative by (4). Notice the source
offset INCREASES the reverse index. This sign is needed to get
the correct reflected index $r-2L$, rather than $r$.

It remains to bound the primitive's $2L$ backward ratios. Put
\[
       R_i=\frac{r+e+3+i}{r-i},\quad0\le i<2L.
\]
All denominators are positive and the skeleton CI envelope gives
$d_r\le(\prod_{i=0}^{2L-1}R_i)d_{r-2L}$.
Pair $i$ with $2L-1-i$. The following elementary inequality holds
for ALL $e,L\ge0$ and $r\ge3e+6L+5$:
\[
       (r+e+3)(r+e+2L+2)\le2r(r-2L+1).               \tag{7}
\]
For example, put $r=3e+6L+5+x$, $x\ge0$; the difference
between right and left is a polynomial with nonnegative
coefficients and positive constant. For each pair let
$u=i(2L-1-i)\ge0$. Its numerator is the left side of (7)
plus $u$, and its denominator is $r(r-2L+1)+u$.
Consequently
\[
             R_iR_{2L-1-i}\le2,
             \qquad\prod_{i=0}^{2L-1}R_i\le2^L.       \tag{8}
\]
This pairing covers every factor exactly once. Combining (4)
and (8) with the primitive curvature bound gives
\[
       C\le\frac{e+1}{r+1}\,2^L d_{r-2L}
          \le2^{L-2}d_{r-2L}.
\]
Equation (6) pays the entire curvature, contradicting $T>0$.
Thus positive recovery requires integer $D\le12d-25$.
Since $h\le D$, $a=\alpha+D+h\le\alpha+24d-50$.
Substituting $d=\kappa-\alpha$ and $\alpha\ge1$ proves
$a\le24\kappa-73$. Also $h\le D\le12d-25\le12\kappa-37$.

For an edge not already paid by the count-mass criterion,
$2g<5j+4$ and $j\le a+1\le24\kappa-72$ imply
$2g<120\kappa-356$. Both sides are even integers after the
last bound, so $g\le60\kappa-179$.
This proves every assertion.
\end{proof}

Nine additional standard-axiom Lean reports check the unbounded
mirror-ratio polynomial, each ratio pair, the COMPLETE arbitrary-length
product pairing, degree-to-rank implications, both scalar payments,
dyadic scaling and the final linear integer bounds. Existing actual
skeleton envelopes, reflection and source-curvature identifications
are joined in the written proof above. Separate finite diagnostics
use exact binomial product ratios and independent actual forward
coefficients; none substitutes for the unbounded argument.

The large-degree region now has pointwise tail nonrecovery, and
therefore all-head S with the usual OLD A2 head input. Small owners
inside the explicit bounds may have genuine positive recoveries;
their head compensation still has to be proved in general.
The opposite-source-prefix premise and the mixed-factor premise
remain essential to this argument. Both-rising late overlap, source
prefixes after decline, general C/S/A2 and original undilated
coefficient unimodality remain OPEN. No original least-counterexample
bound is raised solely from these scoped S refinements.
